% Generated draft: requires mathematical and source review.
\documentclass[11pt]{article}
\usepackage[margin=25mm]{geometry}
\usepackage{fontspec}
\setmainfont{Latin Modern Roman}
\usepackage{amsmath,amssymb,mathtools,unicode-math}
\setmathfont{Latin Modern Math}
\usepackage{enumitem,needspace,xurl,hyperref}
\setlength{\parindent}{0pt}
\setlength{\parskip}{5pt}
\setlength{\emergencystretch}{3em}
\newcommand{\N}{\mathbb{N}}
\newcommand{\Z}{\mathbb{Z}}
\newcommand{\Q}{\mathbb{Q}}
\newcommand{\R}{\mathbb{R}}
\newcommand{\C}{\mathbb{C}}
\newcommand{\F}{\mathbb{F}}
\newcommand{\A}{\mathbb{A}}
\newcommand{\PP}{\mathbb{P}}
\newcommand{\E}{\mathbb{E}}
\newcommand{\eps}{\varepsilon}
\newcommand{\dd}{\,\mathrm{d}}
\begin{document}
\section*{UnsolvedMath ID: 30002898}
\textbf{Source problem code:} catalogue entry
\subsection{Definitions and mathematical statement}
Let \(G\) and \(H\) be countable groups acting by homeomorphisms on compact topological spaces \(X\) and \(Y\), respectively. For the following formulation, give the groups the discrete topology. Call an action topologically free if
\[
\operatorname{int}_X\{x\in X:g\cdot x=x\}=\varnothing
\qquad(g\in G\setminus\{e_G\}),
\]
and impose the analogous condition on \(H\curvearrowright Y\). The source does not specify separation or metrizability assumptions on the spaces, or its precise convention for topological freeness; these remain qualifications requiring review.

Using the continuous-transfer-map convention, the actions are continuously orbit equivalent if there exist a homeomorphism \(h:X\to Y\) and continuous maps
\[
a:G\times X\to H,
\qquad b:H\times Y\to G
\]
such that
\[
h(g\cdot x)=a(g,x)\cdot h(x),
\qquad
h^{-1}(s\cdot y)=b(s,y)\cdot h^{-1}(y)
\]
for every \(g\in G\), \(x\in X\), \(s\in H\), and \(y\in Y\). In particular, \(h(G\cdot x)=H\cdot h(x)\) for every \(x\in X\).

One possible meaning of conjugacy for actions of possibly different groups is the existence of a group isomorphism \(\theta:G\to H\) and a homeomorphism \(k:X\to Y\) satisfying
\[
k(g\cdot x)=\theta(g)\cdot k(x)
\qquad(g\in G,\ x\in X).
\]
The source does not resolve whether this meaning is intended, whether the groups must instead coincide and \(\theta\) must be the identity, or whether the given orbit-equivalence homeomorphism itself must implement conjugacy.

\textbf{Problem, subject to these unresolved conventions.} Characterize the hypotheses on pairs of topologically free actions of countable groups on compact spaces under which continuous orbit equivalence implies conjugacy. With conjugacy interpreted through a group isomorphism as above, the implication to be characterized is
\[
\bigl(\exists\,h,a,b\text{ satisfying the continuous orbit equivalence conditions}\bigr)
\ \Longrightarrow\
\bigl(\exists\,\theta,k\text{ satisfying the conjugacy conditions}\bigr).
\]
In particular, the supplied statement asks for hypotheses beyond a described known duality-group-versus-solvable-group criterion that exclude the described counterexamples for higher-rank free abelian groups and nonabelian free groups. The record supplies neither the exact criterion nor the counterexample constructions or their action hypotheses. Consequently, this refinement cannot be given a precise theorem-level formulation from the supplied evidence alone. The original wording refers to the relevant actions without specifying that class; the displayed formulation records a classification question, not an assertion that the implication holds for all such actions.
\subsection{Short English statement}
Determine which hypotheses force continuously orbit equivalent topologically free actions of countable groups on compact spaces to be conjugate. The supplied record particularly asks for conditions beyond a cited duality-group-versus-solvable-group criterion that exclude described free abelian and nonabelian free group counterexamples. The intended action class, precise meaning of conjugacy, and hypotheses of that criterion require clarification from the cited source.
\subsection{Sources}
\begin{itemize}
\item Local registry: \texttt{data/unsolvedmath/problems.json}, with \texttt{catalog/problems.json} as the additional-record registry; supplied permanent UnsolvedMath ID \(30002898\), title \emph{Rigidity of Continuous Orbit Equivalence}, exact source code \texttt{catalogue entry}. This draft uses the supplied complete record; the local files were not independently inspected.
\item Supplied report citation: \emph{Noncommutative Geometry (2016)}, Oberwolfach Reports, pp.\ 1631--1700, DOI \texttt{10.4171/owr/2015/29}. \url{https://doi.org/10.4171/owr/2015/29}. The record gives proposed year 2015 and does not name an individual proposer. Its background attributes the problem to workshop open-problem discussions collected in Oberwolfach Reports.
\item Supplied literature source: Li, \emph{Continuous orbit equivalence rigidity}. \url{https://arxiv.org/abs/1503.01704}.
\item Supplied literature source: Qiang--Hou, \emph{Acta Mathematica Sinica}. \url{https://doi.org/10.1007/s10114-023-2608-6}.
\item Supplied literature source: \emph{Oberwolfach report 29}. \url{https://doi.org/10.4171/owr/2015/29}.
\item Literature assessment date: \texttt{2026-08-21}. This is supplied dated curation, not independently verified. The record labels the problem partially solved and reports positive and negative rigidity results for important classes, groupoid characterizations of semigroup versions, and no known general reversal of the one-sided implication. It expressly requests independent verification of time-sensitive status and cited claims.
\item The record labels its statement \texttt{corrected\_verified} and reports checking it against the official Oberwolfach report. These are supplied source-verification claims, not verification performed for this draft.
\end{itemize}
\end{document}
